% BHU PYQ -- Manually transcribed & error-corrected
% Paper: BCH-224 : Business Mathematics
% Exam: B.Com. (Hons.) (Semester IV) (Special) Examination, 2014-15

\documentclass[11pt,a4paper]{article}
\usepackage[a4paper,top=2.5cm,bottom=2.5cm,left=2.8cm,right=2.8cm,headheight=14pt]{geometry}
\usepackage{amsmath,amssymb}
\usepackage[T1]{fontenc}
\usepackage[utf8]{inputenc}
\usepackage{lmodern,microtype,fancyhdr,enumitem,hyperref,lastpage,parskip}

\hypersetup{
    colorlinks=true,
    urlcolor=black,
    linkcolor=black,
    pdftitle={BCH-224: Business Mathematics -- BHU B.Com. (Hons.) Sem IV 2014-15}
}

\pagestyle{fancy}
\fancyhf{}
\renewcommand{\headrulewidth}{0.4pt}
\renewcommand{\footrulewidth}{0.4pt}
\fancyhead[L]{\small   \textit{Banaras Hindu University}}
\fancyhead[R]{\small   \textit{BCH-224: Business Mathematics}}
\fancyfoot[C]{\small Page  \thepage\ of \pageref{LastPage}}


\newlist{parts}{enumerate}{2}
\setlist[parts,1]{label=(\alph*),leftmargin=*,itemsep=5pt,topsep=3pt}
\setlist[parts,2]{label=(\roman*),leftmargin=*,itemsep=3pt,topsep=2pt}

\newcommand{\pts}[1]{\hfill   \textit{[#1]}}


\begin{document}


\begin{center}
  {\large\bfseries BANARAS HINDU UNIVERSITY}\\[2pt]
  {\normalsize B.Com. (Hons.) --- Semester IV (Special) Examination, 2014-15}\\[6pt]
  \rule{0.8\linewidth}{0.5pt}\\[6pt]
  {\large\bfseries Paper No. BCH-224: Business Mathematics}\\[4pt]
  {\normalsize Time: 3 Hours\hfill Maximum Marks: 70}
\end{center}

\rule{\linewidth}{0.4pt}

\smallskip



\noindent   \textit{Instructions: Answer all questions. All questions carry equal marks (14 marks each).}

\bigskip




\noindent   \textbf{Question 1.}

\begin{parts}
  \item The 4th term of an A.P. is 13 and the 9th term is 28. Find the 12th term.\pts{7}
  \item The sum of 4 terms of an A.P. is 16 and the sum of the squares of the extreme terms is 50. Find the numbers.\pts{7}
\end{parts}

\centerline{   \textbf{OR}}


\begin{parts}
  \item Find the sum of $n$ odd natural numbers.\pts{7}
  \item If $a, b, c$ are in A.P., prove that $a^2(b + c), b^2(c + a), c^2(a + b)$ are also in A.P.\pts{7}
\end{parts}

\medskip\rule{\linewidth}{0.2pt}




\noindent   \textbf{Question 2.}

\begin{parts}
  \item Which term of the sequence $3, 9, 27, \dots$ is $729$?\pts{7}
  \item Find the Geometric Progression (G.P.) whose first term is 8 and the sum of infinite terms is 16.\pts{7}
\end{parts}

\centerline{   \textbf{OR}}


\begin{parts}
  \item If $a, b, c$ are in G.P. and $x$ and $y$ are the Arithmetic Means between $a$ and $b$, and $b$ and $c$ respectively, prove that:
  \[
  \frac{1}{x} + \frac{1}{y} = \frac{2}{b}
  \]\pts{7}
  \item If $a, b, c$ are in A.P. and $x, y, z$ are in G.P., prove that:
  \[
  x^{b-c} \cdot y^{c-a} \cdot z^{a-b} = 1
  \]\pts{7}
\end{parts}

\medskip\rule{\linewidth}{0.2pt}

\pagebreak




\noindent   \textbf{Question 3.}

\begin{parts}
  \item If ${}^{2n+1}P_{n-1} : {}^{2n-1}P_n = 3:5$, find $n$.\pts{7}
  \item How many numbers can be formed between 10 and 1000 with the help of digits 2, 3, 4, 5, 0, 6, 7, 8, 9? Repetition of digits is not allowed.\pts{7}
\end{parts}

\centerline{   \textbf{OR}}


\begin{parts}
  \item Prove that:
  \[
  {}^{n-1}C_r + {}^{n-1}C_{r-1} = {}^nC_r
  \]\pts{7}
  \item In how many ways can 4 men and 3 ladies be selected from 7 men and 8 ladies?\pts{7}
\end{parts}

\medskip\rule{\linewidth}{0.2pt}




\noindent   \textbf{Question 4.}

\begin{parts}
  \item Find the term independent of $x$ in the expansion of $\left(x^2 + \frac{1}{x}\right)^{12}$.\pts{7}
  \item In the expansion of $(1 + x)^{43}$, the coefficients of $(2r + 1) \text{th}$ term and $(r + 2) \text{th}$ term are equal. Find $r$.\pts{7}
\end{parts}

\centerline{   \textbf{OR}}


\begin{parts}
  \item Find the coefficient of $x^{15}$ in the expansion of $(x - x^2)^{10}$.\pts{7}
  \item Find the 19th term of $\left(2x^{1/2} - y^{1/3}\right)^{20}$.\pts{7}
\end{parts}

\medskip\rule{\linewidth}{0.2pt}




\noindent   \textbf{Question 5.}

\begin{parts}
  \item If $f(x) = \frac{x-1}{x+1}$, prove that $f\left(\frac{x-1}{x+1}\right) = -\frac{1}{x}$.\pts{4}
  \item Evaluate:
  \[
  \lim_{x  o 0} \frac{e^x - e^{-x}}{x}
  \]\pts{4}
  \item Differentiate $y$ w.r.t. $x$:
  \[
  y = \frac{e^{x^3} + e^{-x^3}}{e^{x^3} - e^{-x^3}}
  \]\pts{6}
\end{parts}

\centerline{   \textbf{OR}}


\begin{parts}
  \item Simplify:
  \[
  \int \left(x + \frac{1}{x}\right)^2 \, dx
  \]\pts{4}
  \item Evaluate the following integrals:
  
\begin{parts}
    \item $\int x^3 \, dx$
    \item $\int x^n \log x \, dx$
  \end{parts}\pts{10}
\end{parts}

\vfill
\rule{\linewidth}{0.4pt}

\begin{center}\small
\href{https://bhu-exam.vercel.app/}{Click here to visit our website}\\[4pt]
\href{https://me-aryan.vercel.app/}{Click here to visit our contributor}\\[4pt]
\href{https://quashan-me.vercel.app/}{Click here to visit our contributor}
\end{center}

\end{document}
